\documentclass[10pt]{article}

% COMP 2300 Logic Core companion: student-facing reference and practice.
\usepackage[letterpaper,margin=0.56in,includehead,includefoot,
  headheight=15pt,headsep=7pt,footskip=20pt]{geometry}
\usepackage[T1]{fontenc}
\usepackage{amsmath,amssymb,mathtools}
\usepackage{microtype}
\usepackage[dvipsnames,svgnames]{xcolor}
\usepackage{multicol}
\usepackage{enumitem}
\usepackage{tabularx,booktabs,array}
\usepackage[most]{tcolorbox}
\usepackage{fancyhdr}
\usepackage{titlesec}
\usepackage[hidelinks]{hyperref}

\definecolor{CoreNavy}{HTML}{18364A}
\definecolor{CoreTeal}{HTML}{176B6B}
\definecolor{CoreGold}{HTML}{9A6500}
\definecolor{CoreRed}{HTML}{963D38}
\definecolor{CoreBlueTint}{HTML}{EDF4F8}
\definecolor{CoreTealTint}{HTML}{EDF8F6}
\definecolor{CoreGoldTint}{HTML}{FFF7E5}
\definecolor{CoreRedTint}{HTML}{FCEFED}

\setlength{\parindent}{0pt}
\setlength{\parskip}{2.5pt}
\setlength{\columnsep}{0.25in}
\setlength{\columnseprule}{0.25pt}
\setlength{\emergencystretch}{1.2em}
\setlength{\abovedisplayskip}{3pt}
\setlength{\belowdisplayskip}{3pt}
\setlist{leftmargin=*,itemsep=1pt,topsep=1.5pt,parsep=0pt,partopsep=0pt}
\raggedbottom

\titleformat{\section}{\Large\bfseries\color{CoreNavy}}{}{0pt}{}
  [\vspace{-3pt}\color{CoreNavy}\rule{\linewidth}{0.7pt}]
\titlespacing*{\section}{0pt}{3pt}{4pt}

\pagestyle{fancy}
\fancyhf{}
\fancyhead[L]{\textbf{COMP 2300}}
\fancyhead[C]{\textbf{Logic Core Companion}}
\fancyhead[R]{Definitions and Worked Examples}
\fancyfoot[L]{Core sections A--E}
\fancyfoot[C]{\thepage}
\fancyfoot[R]{Study reference---not used during a Core}
\renewcommand{\headrulewidth}{0.5pt}
\renewcommand{\footrulewidth}{0.3pt}

\tcbset{enhanced,before skip=2.5pt,after skip=2.5pt,boxrule=0.55pt,arc=1.2mm,
  left=4pt,right=4pt,top=3pt,bottom=3pt,fonttitle=\bfseries}
\newtcolorbox{definitionbox}[1]{title={#1},colback=CoreBlueTint,
  colframe=CoreNavy,colbacktitle=CoreNavy,coltitle=white}
\newtcolorbox{methodbox}[1]{title={#1},colback=CoreTealTint,
  colframe=CoreTeal,colbacktitle=CoreTeal,coltitle=white}
\newtcolorbox{workedbox}[1]{title={#1},colback=white,
  colframe=CoreNavy,colbacktitle=CoreBlueTint,coltitle=CoreNavy}
\newtcolorbox{warningbox}[1]{title={#1},colback=CoreRedTint,
  colframe=CoreRed,colbacktitle=CoreRed,coltitle=white}
\newtcolorbox{checklistbox}[1]{title={#1},colback=CoreGoldTint,
  colframe=CoreGold,colbacktitle=CoreGold,coltitle=white}
\newcommand{\term}[1]{\textbf{\color{CoreNavy}#1}}

\begin{document}
\small

\section*{A. Truth Tables and Connectives}
\begin{multicols}{2}
\begin{definitionbox}{Core vocabulary}
A \term{proposition} has a truth value. A \term{valuation} assigns T or F to
each atomic proposition. A \term{tautology} is true on every row, a
\term{contradiction} on none, and a \term{contingency} on some but not all.
Two expressions are \term{logically equivalent} when their final columns match
on every valuation.
\end{definitionbox}

\begin{definitionbox}{Connective meanings}
\[
\begin{array}{c|c|c|c|c}
p&q&\neg p&p\land q&p\lor q\\\hline
T&T&F&T&T\\T&F&F&F&T\\F&T&T&F&T\\F&F&T&F&F
\end{array}
\]
$p\to q$ is false only when $p$ is T and $q$ is F. $p\leftrightarrow q$
is true exactly when $p$ and $q$ match. ``$p$ only if $q$'' and ``$q$ is
necessary for $p$'' both mean $p\to q$.
\end{definitionbox}

\begin{methodbox}{Building an eight-row table}
For atoms $p,q,r$, list rows in a fixed binary pattern:
\[
\begin{array}{c|c|c}
p&q&r\\\hline
T&T&T\\T&T&F\\T&F&T\\T&F&F\\
F&T&T\\F&T&F\\F&F&T\\F&F&F
\end{array}
\]
Add one column for each useful subexpression, working from the innermost
connectives outward. Classify only the final column.
\end{methodbox}


\begin{workedbox}{A.3: a complete truth table}
\textbf{Problem.} Construct a complete eight-row truth table for
$\Phi=(p\to q)\lor(\neg p\to r)$. Show useful intermediate columns and
classify $\Phi$ as a tautology, contradiction, or contingency.

\textbf{Worked solution.} List every valuation once. Compute each implication
before applying the final $\lor$:
\[
\begin{array}{ccc|c|c|c|c}
p&q&r&\neg p&p\to q&\neg p\to r&\Phi\\\hline
T&T&T&F&T&T&T\\
T&T&F&F&T&T&T\\
T&F&T&F&F&T&T\\
T&F&F&F&F&T&T\\
F&T&T&T&T&T&T\\
F&T&F&T&T&F&T\\
F&F&T&T&T&T&T\\
F&F&F&T&T&F&T
\end{array}
\]
As an independent check,
\[
(p\to q)\lor(\neg p\to r)
\equiv(\neg p\lor q)\lor(p\lor r).
\]
Reordering exposes $\neg p\lor p$, which is always true. This checks the
complete table; it does not replace a requested table.

\textbf{Answer.} Every entry in the final column is T, so $\Phi$ is a
\term{tautology}.
\end{workedbox}

\begin{workedbox}{Similar example: interpreting one row}
\textbf{Problem.} Let $p=F$, $q=T$, and $r=F$. Evaluate
$(p\to q)\land(\neg q\lor r)$, showing the intermediate values.

\textbf{Worked solution.}
\[
\begin{aligned}
p\to q&=F\to T=T,\\
\neg q\lor r&=F\lor F=F,\\
(p\to q)\land(\neg q\lor r)&=T\land F=F.
\end{aligned}
\]
Compute the inner pieces first and the main connective last. That dependency
order should match a truth table's columns.

\textbf{Answer.} The proposition is false on this valuation.
\end{workedbox}

\begin{warningbox}{Common errors}
Do not read $p\to q$ as a causal claim or as $q\to p$. Do not omit rows,
change row order midway, or classify from a few sample rows. In
$\neg(p\land q)$ the negation applies to the entire parenthesized expression.
\end{warningbox}

\begin{checklistbox}{A readiness check}
Can you make all eight rows; compute $\neg,\land,\lor,\to,\leftrightarrow$;
show intermediate columns; and classify the final proposition correctly?
\end{checklistbox}
\end{multicols}

\clearpage
\section*{B--C. Translation, Quantifiers, and Negation}
\begin{multicols}{2}
\begin{definitionbox}{Conditionals in ordinary language}
\begin{tabularx}{\linewidth}{@{}lX@{}}
``$p$ if $q$'' & $q\to p$\\
``$p$ only if $q$'' & $p\to q$\\
``$p$ iff $q$'' & $p\leftrightarrow q$\\
``$p$ unless $q$'' & $\neg q\to p$, equivalently $p\lor q$\\
``$p$ whenever $q$'' & $q\to p$
\end{tabularx}
The antecedent is sufficient; the consequent is necessary.
\end{definitionbox}

\begin{workedbox}{Similar example: conditional translation}
\textbf{Problem.} Let $e$ mean ``the account may edit'' and $a$ mean ``the
account is an administrator.'' Translate ``The account may edit only if it is
an administrator'' into propositional logic. Explain the direction of the arrow.

\textbf{Worked solution.} Editing is the claim, and administrator status is a
necessary condition. If editing occurs, administrator status must hold. Thus
$e$ is the antecedent and $a$ is the consequent. Writing $a\to e$ instead
would say every administrator may edit, which the sentence does not promise.

\textbf{Answer.} $e\to a$.
\end{workedbox}

\begin{definitionbox}{Predicates and quantifiers}
A \term{predicate} becomes a proposition after its variables are assigned or
quantified. State the domain. Then
\begin{align*}
\forall x\,P(x)&:\ P\text{ holds for every }x,\\
\exists x\,P(x)&:\ P\text{ holds for at least one }x.
\end{align*}
The order matters: $\exists x\forall y\,R(x,y)$ uses one $x$ for every $y$;
$\forall y\exists x\,R(x,y)$ may use a different $x$ for each $y$.
\end{definitionbox}

\begin{workedbox}{Similar example: quantifiers and argument order}
\textbf{Problem.} Both variables range over the students in a class. Let
$M(x,y)$ mean ``student $x$ messaged student $y$.'' Translate each formula into
English and explain how its meaning differs from the first:
\[
\exists x\forall y\,M(x,y),\quad
\forall y\exists x\,M(x,y),\quad
\exists x\forall y\,M(y,x).
\]

\textbf{Worked solution.} In the first formula, choose $x$ once, before
checking every $y$. In the second, choose a sender $x$ after each recipient
$y$ is specified; the sender may change. In the third, the quantifier order is
the same as the first but the predicate's arguments reverse sender and recipient.

\textbf{Answers.}
\begin{enumerate}
\item One student messaged every student.
\item Every student received a message from at least one student.
\item One student received a message from every student.
\end{enumerate}
Check both quantifier order and predicate argument order. These statements
allow a student to message themself; excluding that requires an explicit condition.
\end{workedbox}


\begin{definitionbox}{Restricting an all-things domain}
If the domain is all things, use predicates to restrict types:
\[
\begin{aligned}
\text{Every }A\text{ is }B &: \forall x(A(x)\to B(x)),\\
\text{Some }A\text{ is }B &: \exists x(A(x)\land B(x)),\\
\text{No }A\text{ is }B &: \forall x(A(x)\to\neg B(x)).
\end{aligned}
\]
Use implication with a universal restriction and conjunction with an
existential restriction.
\end{definitionbox}

\begin{methodbox}{C: push a negation inward}
\[
\neg\forall x\,P\equiv\exists x\,\neg P,
\qquad
\neg\exists x\,P\equiv\forall x\,\neg P.
\]
Work from left to right: (1) flip the first quantifier crossed by the negation,
(2) keep the remaining quantified formula in the same order, and (3) repeat.
Then use De Morgan's laws and conditional elimination until each negation is
immediately before a predicate.
\end{methodbox}

\begin{workedbox}{C.1(a): a complete negation}
\textbf{Problem.} The domain is all things. Define predicates, symbolize
``All dogs have fleas,'' negate it with no negation before a quantifier, and
translate the negation into ordinary English.

\textbf{Worked solution.} Let $D(x)$ mean ``$x$ is a dog'' and $F(x)$ mean
``$x$ has fleas.'' The universal claim restricts attention to dogs with an
implication:
\[
\begin{aligned}
&\forall x(D(x)\to F(x));\\
&\neg\forall x(D(x)\to F(x))\\
&\quad\equiv\exists x\neg(D(x)\to F(x))\\
&\quad\equiv\exists x(D(x)\land\neg F(x)).
\end{aligned}
\]
The last step uses $\neg(P\to Q)\equiv P\land\neg Q$. A counterexample must
be a dog \emph{and} lack fleas, so the existential restriction uses conjunction.

\textbf{Answer.} $\exists x(D(x)\land\neg F(x))$: some dog does not have fleas.
\end{workedbox}

\begin{workedbox}{C.2(b): nested negation}
\textbf{Problem.} Let $P$ be a predicate on two variables over declared,
nonempty domains $X$ and $Y$. Rewrite $\neg\forall x\exists y\,P(x,y)$ so
every negation appears immediately before a predicate. Here $x\in X$ and
$y\in Y$ throughout.

\textbf{Worked solution.}
\[
\begin{aligned}
\neg\forall x\exists y\,P(x,y)
&\equiv\exists x\neg\exists y\,P(x,y)\\
&\equiv\exists x\forall y\neg P(x,y).
\end{aligned}
\]
The first step flips $\forall x$ to $\exists x$; the second flips
$\exists y$ to $\forall y$. Variables remain in their original order.

\textbf{Answer.} $\exists x\forall y\neg P(x,y)$: some $x$ is unrelated by
$P$ to every $y$.
\end{workedbox}

\begin{warningbox}{Common errors}
Do not swap $\forall$ and $\exists$ without moving a negation through them.
Do not reuse a single existential witness unless the formula requires one.
Define both predicate meanings and argument order.
\end{warningbox}

\begin{checklistbox}{B--C readiness}
Can you translate ``if,'' ``only if,'' and ``unless''; distinguish nested
quantifier orders; diagnose reversed arguments; symbolize, negate, simplify,
and translate back to English?
\end{checklistbox}
\end{multicols}

\bigskip
\section*{D. Consistency and Satisfiability}
\begin{multicols}{2}
\begin{definitionbox}{What the question asks}
A collection of formulas is \term{consistent} or \term{satisfiable} when at
least one truth assignment makes all formulas true. It is
\term{inconsistent} when no such assignment exists. A satisfying assignment
is a witness; a contradiction is a proof that no witness exists.
\end{definitionbox}

\begin{methodbox}{A reliable workflow}
\begin{enumerate}
\item Define one atom for each simple statement.
\item Translate every sentence before reasoning.
\item Start with facts and implications whose antecedents are forced.
\item Propagate consequences using valid inference rules.
\item If no contradiction appears, finish a complete assignment and check
every condition.
\item If a contradiction appears, cite a small subset that forces both $p$
and $\neg p$.
\end{enumerate}
\end{methodbox}

\begin{definitionbox}{Useful inference rules}
\begin{align*}
p,\ p\to q&\therefore q &&\text{Modus Ponens}\tag{MP}\\
\neg q,\ p\to q&\therefore\neg p &&\text{Modus Tollens}\tag{MT}\\
p\to q,\ q\to r&\therefore p\to r &&\text{Hypothetical Syllogism}\\
p\lor q,\ \neg p&\therefore q &&\text{Disjunctive Syllogism}
\end{align*}
From $p\leftrightarrow q$, use both $p\to q$ and $q\to p$.
\end{definitionbox}

\begin{workedbox}{D.1: inconsistency and minimality}
\textbf{Problem.} The system is in multiuser state iff it operates normally.
Normal operation implies a functioning kernel. The kernel is not functioning
or the system is in interrupt mode. If the system is not in multiuser state,
it is in interrupt mode. The system is not in interrupt mode. Define atoms,
translate all five conditions, decide consistency, and identify a minimal
inconsistent collection if one exists.

\textbf{Worked solution.} Let $m$ mean multiuser, $n$ normal operation,
$k$ kernel functioning, and $i$ interrupt mode. The conditions are
\[
m\leftrightarrow n,\ n\to k,\ \neg k\lor i,\ \neg m\to i,\ \neg i.
\]
\begin{enumerate}
\item $\neg k$ follows from $\neg k\lor i$ and $\neg i$ (DS).
\item $\neg n$ follows from $n\to k$ and $\neg k$ (MT).
\item $\neg m$ follows from $m\to n$ and $\neg n$ (MT).
\item $i$ follows from $\neg m\to i$ and $\neg m$ (MP), contradicting $\neg i$.
\end{enumerate}
All five conditions are needed. The following valuations satisfy the other
four conditions when the indicated one is removed:
\[
\begin{array}{l|cccc}
\text{Removed condition}&m&n&k&i\\\hline
m\leftrightarrow n&T&F&F&F\\
n\to k&T&T&F&F\\
\neg k\lor i&T&T&T&F\\
\neg m\to i&F&F&F&F\\
\neg i&T&T&T&T
\end{array}
\]
\textbf{Answer.} The conditions are inconsistent; the full five-condition
collection is minimal because deleting any one makes it satisfiable.
\end{workedbox}


\begin{workedbox}{Similar example: one satisfying witness}
\textbf{Problem.} Whenever software is being upgraded, users cannot access
the file system. If users can access it, they can save new files. If users
cannot save new files, the software is not being upgraded. Define atoms,
translate the conditions, and give \emph{one} checked satisfying assignment.

\textbf{Worked solution.} Let $u$ mean software upgraded, $a$ users have
access, and $s$ users can save. The translations are
\[
u\to\neg a,\qquad a\to s,\qquad\neg s\to\neg u.
\]
Try $(u,a,s)=(T,F,T)$. The three formulas evaluate to
\[
T\to T=T,\quad F\to T=T,\quad F\to F=T.
\]
\textbf{Answer.} The collection is consistent, witnessed by $(T,F,T)$.
This task asks for one model. If asked to list \emph{all} models, test all
eight rows or give a complete case split; one witness would be insufficient.
\end{workedbox}

\begin{workedbox}{Similar example: a minimal contradiction}
\textbf{Problem.} Decide whether $p\to q$, $q\to r$, $p$, and $\neg r$ are
consistent. Identify a minimal inconsistent subset and justify minimality.

\textbf{Worked solution.} From $p$ and $p\to q$, infer $q$ by MP. From
$q$ and $q\to r$, infer $r$ by MP, contradicting $\neg r$. To check that
every premise is needed, remove one at a time:
\[
\begin{array}{l|ccc}
\text{Removed}&p&q&r\\\hline
p\to q&T&F&F\\
q\to r&T&T&F\\
p&F&F&F\\
\neg r&T&T&T
\end{array}
\]
Each row satisfies all remaining premises.

\textbf{Answer.} The four-premise set is inconsistent and minimal. A larger
list containing these premises would prove inconsistency but would not itself
answer a request for a minimal collection.
\end{workedbox}

\begin{methodbox}{Listing all models}
For three atoms, a truth table has eight rows. Translate each condition into a
column, conjoin the condition columns, and keep exactly the rows where the
conjunction is T. When reasoning by cases, state why each branch is complete.
\end{methodbox}

\begin{warningbox}{Invalid shortcuts}
Affirming the consequent ($p\to q,q\therefore p$) and denying the antecedent
($p\to q,\neg p\therefore\neg q$) are invalid. ``Only if'' points toward the
necessary condition. A conditional with false antecedent is true.
\end{warningbox}

\begin{checklistbox}{D readiness}
Can you define atoms, translate every condition, give a checked satisfying
assignment or a short contradiction, list all models when asked, identify
forced values, and isolate a minimal contradictory subset?
\end{checklistbox}
\end{multicols}

\clearpage
\section*{E. Equivalence Derivations with Cited Rules}
\begin{multicols}{2}
\begin{definitionbox}{Permitted rules}
Each may be used in either direction:
\begin{align*}
a\to b&\equiv\neg a\lor b &&\text{Implication}\tag{Imp}\\
\neg\neg a&\equiv a &&\text{Double Negation}\tag{DN}\\
\neg(a\lor b)&\equiv\neg a\land\neg b &&\text{De Morgan}\tag{DM}\\
\neg(a\land b)&\equiv\neg a\lor\neg b &&\text{De Morgan}\tag{DM}\\
a\lor(b\land c)&\equiv(a\lor b)\land(a\lor c) &&\text{Distributive}\tag{Dist}\\
a\land(b\lor c)&\equiv(a\land b)\lor(a\land c) &&\text{Distributive}\tag{Dist}\\
(a\circ b)\circ c&\equiv a\circ(b\circ c) &&\text{Associative}\tag{Assoc}\\
a\circ b&\equiv b\circ a &&\text{Commutative}\tag{Comm}\\
a\circ a&\equiv a &&\text{Idempotent}\tag{Idem}
\end{align*}
Here $\circ$ is either $\land$ or $\lor$ where appropriate. Reparenthesizing
requires Associative; reordering requires Commutative.
\end{definitionbox}

\begin{methodbox}{How to plan the chain}
Eliminate implications first. Compare the connective shape of the two sides.
Use Distributive to factor or expand, De Morgan plus Double Negation to move a
negation, and then clean up order or duplicates. Change only one justified
feature per line and cite exactly one rule.
\end{methodbox}

\begin{workedbox}{E.1: factor a common term}
\textbf{Problem.} Show $(p\to q)\land(p\to r)\equiv p\to(q\land r)$. Cite exactly one permitted rule per step.

\textbf{Worked solution.} Rule abbreviations are listed above.
\begin{align*}
(p\to q)\land(p\to r)
&\equiv(\neg p\lor q)\land(p\to r) &&\text{Imp}\\
&\equiv(\neg p\lor q)\land(\neg p\lor r) &&\text{Imp}\\
&\equiv\neg p\lor(q\land r) &&\text{Dist}\\
&\equiv p\to(q\land r) &&\text{Imp}.
\end{align*}
The common disjunct $\neg p$ is what makes the Distributive step possible.

\textbf{Conclusion.} The expressions are equivalent; factoring the common $\neg p$ is the key step.
\end{workedbox}


\begin{workedbox}{E.2: De Morgan structure}
\textbf{Problem.} Show $(p\to r)\land(q\to r)\equiv(p\lor q)\to r$. Cite exactly one permitted rule per step.

\textbf{Worked solution.} Rule abbreviations are listed above.
\begin{align*}
(p\to r)\land(q\to r)
&\equiv(\neg p\lor r)\land(q\to r) &&\text{Imp}\\
&\equiv(\neg p\lor r)\land(\neg q\lor r) &&\text{Imp}\\
&\equiv(r\lor\neg p)\land(\neg q\lor r) &&\text{Comm}\\
&\equiv(r\lor\neg p)\land(r\lor\neg q) &&\text{Comm}\\
&\equiv r\lor(\neg p\land\neg q) &&\text{Dist}\\
&\equiv(\neg p\land\neg q)\lor r &&\text{Comm}\\
&\equiv\neg(p\lor q)\lor r &&\text{DM}\\
&\equiv(p\lor q)\to r &&\text{Imp}.
\end{align*}
The two Commutative steps place the common $r$ where Distributive can factor it.

\textbf{Conclusion.} The expressions are equivalent. The common $r$ is factored before applying De Morgan.
\end{workedbox}

\begin{workedbox}{E.5: nested implication}
\textbf{Problem.} Show $\neg p\to(q\to r)\equiv q\to(p\lor r)$. Cite exactly one permitted rule per step.

\textbf{Worked solution.} Rule abbreviations are listed above.
\begin{align*}
\neg p\to(q\to r)
&\equiv\neg\neg p\lor(q\to r) &&\text{Imp}\\
&\equiv p\lor(q\to r) &&\text{DN}\\
&\equiv p\lor(\neg q\lor r) &&\text{Imp}\\
&\equiv(p\lor\neg q)\lor r &&\text{Assoc}\\
&\equiv(\neg q\lor p)\lor r &&\text{Comm}\\
&\equiv\neg q\lor(p\lor r) &&\text{Assoc}\\
&\equiv q\to(p\lor r) &&\text{Imp}.
\end{align*}
Each line changes exactly one feature and cites exactly one permitted rule.

\textbf{Conclusion.} The expressions are equivalent, with one cited permitted rule at each step.
\end{workedbox}

\begin{warningbox}{Common errors}
Do not cite ``algebra'' or ``simplify.'' Do not silently move parentheses or
terms. Do not use a truth table when a rule-by-rule derivation is requested.
Make sure every intermediate expression is equivalent to the preceding one.
\end{warningbox}

\begin{checklistbox}{Whole-Core readiness}
Can you complete a truth table; translate conditional and quantified language;
negate nested quantifiers; prove consistency or inconsistency; and write an
equivalence chain with one named rule per step? A live form selects one problem
from each section A--E and is designed for 20 minutes.
\end{checklistbox}
\end{multicols}

\end{document}
